
\documentclass{article}
\usepackage{colortbl}
\usepackage{makecell}
\usepackage{multirow}
\usepackage{supertabular}

\begin{document}

\newcounter{utterance}

\centering \large Interaction Transcript for game `wordle', experiment `medium\_frequency\_words\_no\_clue\_no\_critic', episode 5 with reflexion{-}Qwen3.5{-}9B.
\vspace{24pt}

{ \footnotesize  \setcounter{utterance}{1}
\setlength{\tabcolsep}{0pt}
\begin{supertabular}{c@{$\;$}|p{.15\linewidth}@{}p{.15\linewidth}p{.15\linewidth}p{.15\linewidth}p{.15\linewidth}p{.15\linewidth}}
    \# & \multicolumn{2}{c}{Player} && \multicolumn{2}{c}{Game Master} \\
    \hline

    \theutterance \stepcounter{utterance}  
    & & & \multicolumn{4}{p{0.6\linewidth}}{
        \cellcolor[rgb]{0.9,0.9,0.9}{
            \makecell[{{p{\linewidth}}}]{
                \texttt{\tiny{[P1$\langle$GM]}}
                \texttt{You have attempted this same task before and did not succeed. These are the plans you wrote after those attempts.} \\
\texttt{Trial \#1: New plan: The last attempt was aborted because "ates" is not a valid five{-}letter English word (it is a suffix or plural of "at", but not a standalone word in standard Wordle dictionaries). Additionally, my explanation incorrectly claimed 'n' was incorrect when the previous feedback for "atone" showed 'n' as red, which is consistent, but I failed to construct a valid word with the confirmed letters 't' (index 1), 'a' (misplaced), and 'e' (misplaced). Next time, I will strictly ensure my guess is a valid dictionary word before submitting. Given 't' is at index 1, 'a' and 'e' are present but not at indices 0 and 4 respectively, and 'o', 'n' are out, I should try a word like "tears" or "teams" to test the remaining positions while maintaining validity. I will also double{-}check the word list before finalizing the guess.} \\
\\ 
\\ 
\texttt{You are a language wizard who likes to guess words by using the given rules.} \\
\\ 
\texttt{Welcome to Wordle! You have six attempts to guess the target word, a valid English word of five lowercase letters (a{-}z). Please use the tags "explanation:" and "guess:" to provide a concise explanation for each guess.} \\
\\ 
\texttt{For instance, if your guess is "apple", your response should be} \\
\texttt{explanation: this is a common five{-}letter English word, and I am starting my guess with this word.} \\
\texttt{guess: apple} \\
\\ 
\texttt{After each guess, your answer will be validated, and you will receive feedback indicating which letters are correct (green), which letters are correct but in the wrong position (yellow), and which letters are incorrect (red). This feedback can be useful in determining which letters to include or exclude in your next guess.} \\
\\ 
\texttt{For example, the feedback for "apple" might be:} \\
\texttt{guess\_feedback: a<yellow> p<yellow> p<green> l<yellow> e<red>} \\
\\ 
\texttt{The explanation should contain details about how the guess\_feedback is used to arrive at a new guess.} \\
\\ 
\texttt{Do not generate any other text.} \\
\\ 
\texttt{Let's begin with your first guess.} \\
            }
        }
    }
     \\ \\

    \theutterance \stepcounter{utterance}  
    & \multicolumn{4}{p{0.6\linewidth}}{
        \cellcolor[rgb]{0.561,0.737,0.561}{
            \makecell[{{p{\linewidth}}}]{
                \texttt{\tiny{[P1$\rangle$GM]}}
                \texttt{explanation: I will start with "crane" as it contains common vowels and consonants to maximize information gain in the first guess.} \\
\texttt{guess: crane} \\
            }
        }
    }
    & & \\ \\

    \theutterance \stepcounter{utterance}  
    & & & \multicolumn{2}{p{0.3\linewidth}}{
        \cellcolor[rgb]{0.9,0.9,0.9}{
            \makecell[{{p{\linewidth}}}]{
                \texttt{\tiny{[GM$|$GM]}}
                \texttt{attempts = 1} \\
\texttt{target\_word = steam} \\
\texttt{guess = crane} \\
\texttt{guess\_feedback = c<red> r<red> a<yellow> n<red> e<yellow>} \\
            }
        }
    }
    & & \\ \\

    \theutterance \stepcounter{utterance}  
    & & & \multicolumn{4}{p{0.6\linewidth}}{
        \cellcolor[rgb]{0.9,0.9,0.9}{
            \makecell[{{p{\linewidth}}}]{
                \texttt{\tiny{[P1$\langle$GM]}}
                \texttt{guess\_feedback: c<red> r<red> a<yellow> n<red> e<yellow>} \\
\\ 
\texttt{Provide your response only in this format.} \\
\texttt{explanation: details} \\
\texttt{guess: word} \\
            }
        }
    }
     \\ \\

    \theutterance \stepcounter{utterance}  
    & \multicolumn{4}{p{0.6\linewidth}}{
        \cellcolor[rgb]{0.561,0.737,0.561}{
            \makecell[{{p{\linewidth}}}]{
                \texttt{\tiny{[P1$\rangle$GM]}}
                \texttt{explanation: From the feedback on "crane", I know 'a' and 'e' are in the word but not at indices 2 and 4 respectively. 'c', 'r', and 'n' are incorrect. I need to find a word with 'a' and 'e' in different positions. Since 'a' is not at index 2 and 'e' is not at index 4, I will try placing 'a' at index 0 and 'e' at index 3, while using common letters like 't' and 's' to fill the gaps. "tears" fits this pattern: 't' at 0, 'e' at 1 (checking if 'e' can be elsewhere), 'a' at 2 (checking if 'a' can be elsewhere), 'r' at 3 (red), 's' at 4. Wait, 'r' was red in "crane", so 'r' is out. Let's try "teams": 't' at 0, 'e' at 1, 'a' at 2, 'm' at 3, 's' at 4. But 'a' was yellow at index 2 in "crane", meaning 'a' is correct but not at index 2. So "teams" puts 'a' at index 2, which is wrong. I need 'a' not at 2. How about "teats"? 't' at 0, 'e' at 1, 'a' at 2 (wrong position), 't' at 3, 's' at 4. Still 'a' at 2. Let's try "steal": 's' at 0, 't' at 1, 'e' at 2 (wrong, 'e' was yellow at 4), 'a' at 3 (wrong, 'a' was yellow at 2), 'l' at 4. This tests if 'e' is at 2 and 'a' is at 3. But I know 'a' is not at 2 and 'e' is not at 4. Let's try "state": 's' at 0, 't' at 1, 'a' at 2 (wrong), 't' at 3, 'e' at 4 (wrong). No. How about "seats"? 's' at 0, 'e' at 1, 'a' at 2 (wrong), 't' at 3, 's' at 4. Still 'a' at 2. I need a word where 'a' is not at 2 and 'e' is not at 4. Let's try "teach": 't' at 0, 'e' at 1,} \\
            }
        }
    }
    & & \\ \\

    \theutterance \stepcounter{utterance}  
    & & & \multicolumn{2}{p{0.3\linewidth}}{
        \cellcolor[rgb]{0.9,0.9,0.9}{
            \makecell[{{p{\linewidth}}}]{
                \texttt{\tiny{[GM$|$GM]}}
                \texttt{Error: The response should contain the 'guess:' keyword exactly once.} \\
            }
        }
    }
    & & \\ \\

    \theutterance \stepcounter{utterance}  
    & & & \multicolumn{2}{p{0.3\linewidth}}{
        \cellcolor[rgb]{0.9,0.9,0.9}{
            \makecell[{{p{\linewidth}}}]{
                \texttt{\tiny{[GM$|$GM]}}
                \texttt{game\_result = ABORT} \\
            }
        }
    }
    & & \\ \\

    \theutterance \stepcounter{utterance}  
    & & \multicolumn{4}{p{0.6\linewidth}}{
        \cellcolor[rgb]{0.95,0.95,0.95}{
            \makecell[{{p{\linewidth}}}]{
                \texttt{\tiny{[P1$\leftrightarrow$P]}}
                \texttt{Post{-}episode Reflexion (not sent to GM):} \\
\texttt{New plan: The last attempt was aborted because I generated an incomplete word ("teach" cut off mid{-}sentence) and failed to finalize the response within the required format before the turn ended. I also got lost in a loop of hypothetical scenarios without committing to a specific, valid five{-}letter word. Next time, I will immediately select a complete, valid five{-}letter word that adheres to the constraints derived from the previous feedback (e.g., avoiding 'c', 'r', 'n'; placing 'a' and 'e' in new positions). I will construct the full explanation and the complete guess string in one go, ensuring no text is truncated, and strictly follow the "explanation:" and "guess:" tags without adding any trailing thoughts or incomplete sentences.} \\
            }
        }
    }
    & \\ \\

\end{supertabular}
}

\end{document}
